% P-F65 CANDIDATE RESEARCH NOTE. Not yet promoted to the checked paper.
% Written during checkpoint08 freeze in a separate new file; full peer checks pending.
\subsection{Candidate P-F65: the second tensor component has a free smooth slice}

This is an unpromoted candidate proof, recorded separately from the checked
paper.  Its conclusions require complete whole-proof review before use.
Let $U_n$ denote the balanced tensor called $K_n$ in P-F64, and use
$n\ge4$, $M=2^{n-1}$, $b_1=0$, and $b_i=2^{n-i}$ for $2\le i\le n$.
Put $D_r=\operatorname{diag}(r^{b_i})$ for $r^M=1$.
The candidate conclusions are
\[
 \operatorname{Stab}_{GL_n^3}(U_n)
 =\{(aD_r^{-1},bD_r,cD_r):abc=1,\ r^M=1\},             \tag{P-F65a}
\]
and the following consequences.  The effective $SL_n^3$ stabilizer is
trivial; the quotient is smooth at $q_n$; its $\mu_M$ fixed tangent
space has dimension one; and the fixed component $Z_n$ through $q_n$
is exactly its grading-orbit closure.  Its normalized leading density
is one.  For $\zeta$ primitive of order $M$, its local determinant
coefficient is
\[
 \mathcal C_n(\zeta)=
 \frac{\displaystyle\prod_{i\ne j}(1-\zeta^{b_i-b_j})^3}
 {\displaystyle\prod_{i,j,k:\ b_i-b_j-b_k\not\equiv0\ (M)}
                  (1-\zeta^{b_i-b_j-b_k})}.             \tag{P-F65b}
\]
This is a local coefficient, not a claim that the full Hilbert series
has a pole or that the sum over all components is nonzero.

\paragraph{1. A simultaneous diagonal normalization.}
Identify the last two factors of $U_n$.  Its symmetric slice quadrics
have the form
\[
 q_1=\beta y_1^2+\gamma_2y_2^2,\qquad
 q_i=2\alpha_i y_1y_i+\gamma_{i+1}y_{i+1}^2\ (2\le i<n),
 \qquad q_n=2\alpha_n y_1y_n,
\]
with all coefficients nonzero.  A diagonal change $y_i\mapsto t_i y_i$
in both last factors, and independent diagonal rescaling of the slice
basis in the first factor, reduce these to
\[
 q_1=y_1^2+y_2^2,\quad
 q_i=2y_1y_i+y_{i+1}^2\ (2\le i<n),\quad q_n=2y_1y_n. \tag{P-F65c}
\]
Indeed choose $t_1\ne0$, then
$t_2^2=(\beta/\gamma_2)t_1^2$ and recursively
$t_{i+1}^2=(\alpha_i/\gamma_{i+1})t_1t_i$; all choices are nonzero
over $\mathbb C$.  Rescale $q_1$ by $(\beta t_1^2)^{-1}$ and
$q_i$ by $(\alpha_i t_1t_i)^{-1}$ for $i\ge2$.
All these changes are diagonal, so they commute with every displayed
$D_r$.  It therefore suffices to determine the stabilizer in the
normal form (P-F65c).

\paragraph{2. The partial centroid is scalar.}
Write $q_i(y)=y^{\mathsf T}M_i y$, with symmetric matrices $M_i$.
Suppose $E M_i=M_iE^{\mathsf T}$ for every $i$.  This implies that
$E$ preserves the image of every $M_i$.  In particular it preserves
\[
 \operatorname{im}M_1=\langle e_1,e_2\rangle,
 \qquad \operatorname{im}M_n=\langle e_1,e_n\rangle,
\]
whose intersection is $\mathbb C e_1$, since $n\ge4$.
Thus $Ee_1=a e_1$.  Symmetry of $EM_1$ then gives $Ee_2=c e_2$.
The first two columns of $EM_2$ are $c e_2,a e_1$, so their
symmetry forces $c=a$; its third column is $Ee_3$, and symmetry
with every other column forces $Ee_3$ to be a scalar multiple of
$e_3$.  Inductively, if $Ee_i=c_i e_i$, then symmetry of $EM_i$
forces $c_i=a$ and $Ee_{i+1}$ to be a scalar multiple of $e_{i+1}$.
Finally symmetry of $EM_n$ forces that last scalar also to be $a$.
Hence $E=aI$.

The tensor is symmetric in its last two factors.  If $(A,B,C)$ fixes
it, then so does $(A,C,B)$, and consequently
$(I,B^{-1}C,C^{-1}B)$ fixes it.  With $E=B^{-1}C$, multiplying the
tensor equality by $I\otimes I\otimes E$ gives
$(I,E,I)U_n=(I,I,E)U_n$.  This is exactly the matrix equation just
analyzed.  Therefore $C$ is a scalar multiple of $B$.
The full product stabilizer reduces to the congruence stabilizer of
the slice space $Q=\langle q_1,\ldots,q_n\rangle$.

\paragraph{3. The complete rank-two locus, including $n=4$.}
The symmetric matrix of $q=\sum_i c_iq_i$ has entries
\[
 M_{11}=c_1,\qquad M_{1j}=M_{j1}=c_j\ (j\ge2),\qquad
 M_{jj}=c_{j-1}\ (j\ge2),
\]
and all other entries zero.  Let $k$ be the number of nonzero
entries among $c_1,\ldots,c_{n-1}$.  Its bottom-right diagonal
block has rank $k$, so rank at most two requires $k\le2$.
If $k>0$ and some $c_j\ne0$ with $j\ge2$ but $c_{j-1}=0$,
eliminating the invertible $k$ by $k$ diagonal subblock leaves
a principal block $\left(\begin{smallmatrix}u&c_j\\c_j&0\end{smallmatrix}\right)$.
It has rank two, so the whole matrix has rank at least $k+2>2$.
Thus in the cases $k=1,2$, every nonzero $c_j$ with $j\ge2$
must have nonzero predecessor.  The nonzero entries form a prefix:
for $k=1$ only $c_1$ survives, giving $[q_1]$; for $k=2$ only
$c_1,c_2$ survive.  Here $c_n=0$ because $n\ge4$, and the
nonzero three-dimensional block is
\[
 \begin{pmatrix}c_1&c_2&0\\c_2&c_1&0\\0&0&c_2\end{pmatrix}.
\]
It has rank two exactly when $c_1^2=c_2^2$, giving
$[q_1+q_2]$ and $[q_1-q_2]$.
For $k=0$, only $c_n$ can be nonzero, giving $[q_n]$.
These exhaust the projective rank-at-most-two locus:
\[
                 [q_1],\ [q_1+q_2],\ [q_1-q_2],\ [q_n]. \tag{P-F65d}
\]
The first three are distinct collinear points; the last is off
that line.  Therefore every congruence preserving $Q$ preserves
$\mathbb Cq_n$ and the plane $\langle q_1,q_2\rangle$.

\paragraph{4. No exchange of the first and last coordinate lines.}
The equation $q_n=2y_1y_n$ shows that its congruence stabilizer
preserves the unordered pair of factor lines
$\{\mathbb Cy_1,\mathbb Cy_n\}$.
The smallest space of linear forms containing all variables of
$\langle q_1,q_2\rangle$ is $\langle y_1,y_2,y_3\rangle$;
equivalently it is the span of their first derivatives.  It too
is preserved.  Since $y_1$ belongs to this space and $y_n$ does
not for $n\ge4$, the two factor lines cannot be exchanged.
Hence a preserving substitution has $y'_1=t_1y_1$ and
$y'_n=t_ny_n$.

Modulo $y_1$, the image of $Q$ is exactly the diagonal quadratic
space $\langle y_2^2,\ldots,y_n^2\rangle$.
An invertible linear substitution preserving that space is
monomial: the square of each transformed coordinate has no
mixed term, so in characteristic zero that coordinate has at
most one nonzero coefficient.  Thus, before lifting back,
\[
 y'_i=t_i y_{\sigma(i)}+a_i y_1\quad(2\le i\le n),
 \qquad \sigma(n)=n,\ a_n=0,
\]
where $\sigma$ is a permutation of $2,\ldots,n$ and $t_i\ne0$.
Proceed backwards from $i=n-1$ to $i=2$.  If
$y'_{i+1}=t_{i+1}y_{i+1}$, then
\[
 q_i(y')=2t_1t_i y_1y_{\sigma(i)}+2t_1a_i y_1^2
                                  +t_{i+1}^2y_{i+1}^2.
\]
The sole pure square forces the coefficient of $q_i$ to be
$t_{i+1}^2$ and all other $q_j$ with $j<n$ to vanish.
There is no $y_2^2$ term, hence no independent $y_1^2$ term.
Only a multiple of $q_n$ may additionally occur.  Since
$\sigma(i)\ne n$, comparison gives
\[
       \sigma(i)=i,\qquad a_i=0,\qquad t_1t_i=t_{i+1}^2.
\]
Finally $q_1(y')\in Q$ forces $t_1^2=t_2^2$.
Writing $r=t_n/t_1$, these equations say
$t_i/t_1=r^{2^{n-i}}$ for $i\ge2$ and $r^M=1$.
The congruence is precisely a scalar times $D_r$.

Under this congruence $q_i(D_r y)=r^{b_i}q_i(y)$ for every $i$.
The first factor that restores the tensor is consequently the
inverse diagonal matrix, up to the independent scalar factors.
The coefficient of $e_1^{\otimes3}$ forces their product to be
one.  This proves (P-F65a) for the normal form and, by the commuting
diagonal normalization, for the original balanced tensor.

\paragraph{5. The effective quotient and the fixed component.}
The determinant-product character of (P-F65a) is
$\det D_r=r^{M-1}=r^{-1}$, since
$\sum_i b_i=M-1$.  It is faithful on $\mu_M$.
In particular an element in $SL_n^3$ must have $r=1$ and consists
only of scalar matrices whose product is one.  This is exactly
the ineffective central kernel.  The balanced orbit is closed
by P-F64, so the effective free-slice application of Luna used
in P-F60 makes the quotient smooth at $q_n$.
The same closed-orbit lifting argument as in P-F60 shows that
the grading stabilizer of $q_n$ is exactly $\mu_M$.

For grading $\zeta$, use the fixing lift
$(D_\zeta,D_\zeta^{-1},D_\zeta^{-1})$, whose character is $\zeta$.
Its tensor weights are $b_i-b_j-b_k$.  The set of distinct
$b_i$ is $P=\{0,1,2,\ldots,2^{n-2}\}$, where after $1$ only
powers of two occur.  There are exactly $3n-2$ zero weights
modulo $M$.  If one of $b_j,b_k$ is zero, their sum is in $P$,
giving $2n-1$ ordered pairs.  If both are nonzero, their sum is
in $P$ modulo $M$ exactly when they are equal, giving $n-1$
more pairs, including the pair whose sum is $M$.
Two distinct positive powers have a two-bit sum strictly less
than $M$ and cannot give another possibility.
The fixed adjoint subspace of each $\mathfrak{sl}_n$ consists
of its diagonal traceless matrices, dimension $n-1$.
Exactness of finite-group invariants and the free orbit tangent
therefore give
\[
             \dim(T_{q_n}X_n)^{\mu_M}
                      =(3n-2)-3(n-1)=1.               \tag{P-F65e}
\]
The fixed scheme is smooth at this smooth point in characteristic
zero, as in the formal averaging argument of P-F23.
Its component through $q_n$ thus has dimension one and is the
grading-orbit closure $Z_n$.

The restriction algebra on that orbit closure is a monomial
curve $\mathbb C[u^s:s\in S]$.  Exact stabilizer $\mu_M$ says
$\gcd S=M$, and finite generation implies $S/M$ is cofinite.
Consequently its Hilbert function is eventually one precisely
on the multiples of $M$, and
\[
                 \lim_{t\to\zeta}(1-t^M)H_{Z_n}(t)=1. \tag{P-F65f}
\]
This establishes its local leading density with the same
normalization used in P-F23 and P-F60.

\paragraph{6. Local determinant and its exact phase.}
The moving tangent representation is the moving tensor
representation minus the three moving adjoint representations.
For the pullback convention $F\mapsto F\circ h$, the cotangent
eigenvalues equal the point differential eigenvalues; this is
a transpose, not the inverse representation convention.
The first adjoint weights are $b_i-b_j$ and the other two are
their negatives.  Products over all ordered $i\ne j$ are equal
after this sign reversal.  This proves (P-F65b), with density
one from (P-F65f).

For the principal embedding $\zeta=\exp(2\pi i/M)$, put
$N=n^3-3n+2$, the number of moving tensor weights; it is even.
The numerator of (P-F65b) is positive real, by pairing the
ordered adjoint weights of opposite signs.  For each moving
weight let $e$ be its representative in $\{1,\ldots,M-1\}$.
Since $1-\zeta^e=-2i\exp(\pi i e/M)\sin(\pi e/M)$,
\[
 \mathcal C_n(\zeta)=(-1)^{N/2+L}
       \exp(-\pi i n^2/M)|\mathcal C_n(\zeta)|,
 \qquad L=\frac{\sum e-n^2}{M}\in\mathbb Z.           \tag{P-F65g}
\]
Here integrality follows also from the determinant character:
the sum of all tensor weights is $-n^2(M-1)$ modulo $M$.
For a direct evaluation of its parity, all raw weights lie
in $[-M,M/2]$, and precisely one equals $-M$.
The number $P_+$ of strictly positive raw weights is
\[
 P_+=\sum_{k=1}^{n-1}k^2-(n-2).                        \tag{P-F65h}
\]
Indeed, below $2^r$ there are $r+1$ allowed values including
zero; all their ordered sums are smaller than $2^r$ except the
pair $(2^{r-1},2^{r-1})$ when $r\ge1$.
The sum of the moving raw weights is $-n^2(M-1)+M$.
There are $N-P_+$ negative moving weights, each of which adds
$M$ on reduction.  Thus
$L=1-n^2+N-P_+$ and
$N/2+L\equiv N/2+1+n^2+P_+\pmod2$.
Substituting (P-F65h), or checking the four residue classes
of $n$ modulo four, gives
\[
 \mathcal C_n(e^{2\pi i/M})=
 \begin{cases}-1,&n\equiv1\pmod4,\\1,&n\not\equiv1\pmod4\end{cases}
 \exp(-\pi i n^2/M)|\mathcal C_n(e^{2\pi i/M})|.        \tag{P-F65i}
\]
For another primitive embedding $\zeta^a$, the always-valid
phase formula instead uses
$L_a=(\sum [a(b_i-b_j-b_k)]_M-a n^2)/M$ on the moving set.
No uniform sign across embeddings or across different fixed
components has been deduced.

The geometric inputs here are exactly the already checked
Kempf--Ness, affine quotient, and free-slice statements cited
in P-F60, and the known balanced support cited precisely in
P-F64.  The stabilizer, rank-locus, fixed-count, and phase
arguments above are supplied in full for independent review.
